% capitoli/02_stato_arte.tex - Capitolo 2: Stato dell'Arte

\chapter{Stato dell'Arte}
\label{ch:stato_arte}

\section{Process Mining in Ambito Clinico}
Il settore sanitario genera quotidianamente moli di dati strutturati e non strutturati che riflettono l'interazione dei pazienti con le strutture ospedaliere, rendendo i \textit{Log degli Eventi} una miniera di informazioni fondamentale per analizzare oggettivamente i percorsi di cura. La gestione e il miglioramento di questi percorsi sono sfide affrontate storicamente attraverso il \textit{Process Mining}.

In ambito clinico, il Process Mining consiste in una serie di tecniche volte a estrarre modelli di processo dai dati (ad esempio tracce esportate in formato XES) per scoprire, monitorare e ottimizzare l'erogazione delle cure mediche e l'efficienza dei servizi. Gli algoritmi di Process Discovery o di Conformance Checking permettono alle amministrazioni di individuare anomalie procedurali e deviazioni dai percorsi terapeutici codificati (Clinical Pathways). Sebbene questi metodi godano di robustezza matematica per mappare il sistema, si sono spesso rivelati rigidi di fronte alle variabili caotiche e all'alta stocasticità dei flussi ospedalieri, che esibiscono dipendenze temporali complesse difficili da cogliere attraverso soli grafi deterministici.

\section{Evoluzione del Deep Learning: Da RNN a Transformer}
La necessità di mitigare le carenze del Process Mining basato su regole si è tradotta nell'integrazione di tecniche di Data Mining esplorativo, giungendo ben presto all'utilizzo del \textit{Deep Learning}. 

Inizialmente, l'esigenza di elaborare tracce cliniche di lunghezza variabile ha portato all'impiego delle architetture ricorrenti, in primis Recurrent Neural Networks (RNN) e Long Short-Term Memory (LSTM), poiché naturalmente predisposte per dati sequenziali. Tuttavia, il difetto cronico delle RNN legato alla dispersione del gradiente (\textit{vanishing gradient}) e l'incapacità intrinseca di gestire parallelismi e interdipendenze a lungo raggio in storie cliniche complesse ne hanno fortemente limitato l'efficacia pratica.

Il cambio di paradigma epocale avviene nel 2017 con l'introduzione dell'architettura Transformer e l'invenzione dei meccanismi di \textit{Self-Attention}. Un passo decisivo verso l'analisi semantica delle tracce ospedaliere è stato il superamento delle classiche codifiche One-Hot, tipiche dell'era pre-Transformer, in favore di \textit{embeddings} linguistici. L'avvento dei Large Language Models (LLM), tra cui Bidirectional Encoder Representations from Transformers (BERT) \cite{devlin2019bert}, ha consentito di addestrare vettori profondamente contestuali su moli sterminate di dati testuali non etichettati, estraendo conoscenze sintattiche e logiche che vengono trasferite con successo a task applicativi specifici. Nel caso clinico, ciò si traduce in un modello capace di osservare e valutare un'intera traccia di eventi contemporaneamente e in modo bidirezionale, anziché passo dopo passo nel tempo.

\section{L'Approccio Storytelling: Il Progetto LEGOLAS}
Affinché i modelli Transformer, ed in particolare BERT, possano elaborare con profitto i dati tratti da log storici di Process Mining, si è resa necessaria una forte innovazione nella tecnica di rappresentazione dell'informazione (\textit{embedding}). Nasce da questa esigenza l'approccio \textit{Storytelling}.

L'idea fondamentale, che ispira anche il presente lavoro di tesi, si radica nel recente e pionieristico progetto \textit{LEGOLAS} (Leveraging a Large Language Model to Predict Hospital Admissions of Emergency Department Patients) \cite{PASQUADIBISCEGLIE2025128224}. L'intuizione principale di LEGOLAS consiste nel superare la mera concatenazione meccanica di eventi e timestamp discreti per spingersi nella generazione di veri e propri documenti di testo coerenti, in linguaggio naturale. 

La traccia di un paziente non è più vissuta dall'intelligenza artificiale come una sequela di array numerici estratti dal registro dell'ospedale, ma viene trascritta in sotto forma di "storia" cronologica, in modo che il testo contenga esplicitamente la semantica degli eventi temporali e delle distanze tra un'azione clinica e l'altra. Poiché BERT è un LLM pre-addestrato sul linguaggio naturale umano, nutrirlo con "racconti clinici" formattati linguisticamente ad hoc permette un salto prestazionale considerevole in fase di Fine Tuning per task di classificazione, mitigando problemi di dimensionalità esplosiva tipici del Process Mining.

\section{Explainable AI (XAI) per i Modelli Transformer}
Per superare i timori sulla natura a scatola nera delle reti neurali, l'Explainable AI (XAI) emerge come necessità non eludibile. Nel contesto del Deep Learning per il settore medico, gli algoritmi trasparenti come il Decision Tree non reggono in termini di performance rispetto al Deep Learning, rendendo cruciale sviluppare metodi post-hoc (che si applicano a modello addestrato) per interpretare reti complesse.

Nel caso dei modelli basati su Transformer, si analizza il focus cognitivo della rete ricorrendo comunemente alle \textit{Mappe di Attenzione} della rete stessa, un proxy del processo di ragionamento interno del modello. Tuttavia, le \textit{Attention Maps} riflettono i "pesi" posizionali generici, che non indicano un contributo deduttivo diretto verso una specifica decisione di classificazione, e sono note per essere affette da \textit{noise} e scarsa proporzionalità reale all'importanza di un token.

Al fine di ottenere metriche di feature importance rigorose rispetto alla classe predetta (es. "degente lungo termine"), la letteratura scientifica più recente suggerisce di guardare al calcolo dei Gradienti Integrati. Proposto originariamente da Sundararajan et al. \cite{sundararajan2017axiomatic}, il metodo degli \textit{Integrated Gradients} (IG) è una tecnica assiomatica e robusta. L'intuizione dell'algoritmo IG consiste nell'accumulare i gradienti lungo una linea retta nello spazio delle feature, passando gradualmente da un input "baseline" neutro fino all'input reale testuale (la storia del paziente), garantendo matematicamente che la somma del contributo di tutti i singoli token corrisponda esattamente allo scarto tra la probabilità emessa sulla baseline neutra e quella finale calcolata. Questa aderenza matematica (\textit{Completeness}) fa degli IG lo strumento d'elezione per misurare minuziosamente quale parola ha spinto o inibito la decisione del classificatore BERT.
